\documentclass[12pt]{article}
\usepackage[margin=0.65in]{geometry}

\newcommand{\duedate}{}
\newcommand{\assignment}{Differentiable Neural Computers}

\newcommand{\name}{Aksel Kretsinger-Walters}
\newcommand{\email}{}

\makeatletter
\def\input@path{{../}{../../}{../../../}}
\makeatother
\input{pset_template_CLMM.tex}

\begin{document}
\psetheader

\section*{Differentiable Neural Computers}

\paragraph{Motivation.}
Standard neural networks struggle with tasks requiring structured memory and algorithmic reasoning.
Recurrent networks can store information in hidden states, but this capacity is limited and difficult
to control. Differentiable Neural Computers (DNCs) augment neural networks with an explicit,
differentiable external memory that allows models to store, retrieve, and manipulate information
similar to traditional computers.

\paragraph{Architecture.}
A DNC consists of two main components:

\begin{itemize}
		\item A \textbf{controller} (typically an LSTM) that processes inputs.
		\item An external \textbf{memory matrix} that stores information across time.
\end{itemize}

The memory is represented as

\[
		M_t \in \mathbb{R}^{N \times W}
\]

where $N$ is the number of memory slots and $W$ is the width of each slot.

At each timestep the controller produces an interface vector which determines:

\begin{itemize}
		\item where to write
		\item where to read
		\item what content to store
\end{itemize}

These operations are differentiable, enabling end-to-end training with gradient descent.

\paragraph{Memory Addressing.}
The model uses two addressing mechanisms.

\textbf{Content-based addressing}

\[
		w^c_i = \frac{\exp(\beta K(k, M_i))}{\sum_j \exp(\beta K(k, M_j))}
\]

where $K(\cdot,\cdot)$ measures similarity between the key and memory rows.

\textbf{Temporal addressing}

The DNC tracks a \emph{temporal link matrix} that records the order in which memory
locations were written. This enables forward and backward traversal through stored
sequences.

\paragraph{Memory Allocation.}
To avoid overwriting useful information, the model tracks memory usage.
Unused locations are allocated preferentially when writing new information.

This enables the system to dynamically grow stored data structures
during computation.

\paragraph{Read and Write Heads.}

The network maintains multiple read and write heads.

Write heads determine

\[
		M_t = M_{t-1} (1 - w^w_t e_t^\top) + w^w_t v_t^\top
\]

where

\begin{itemize}
		\item $w^w_t$ is the write weighting
		\item $e_t$ is the erase vector
		\item $v_t$ is the write vector
\end{itemize}

Read heads produce read vectors

\[
		r_t = M_t^\top w^r_t
\]

which are returned to the controller.

\paragraph{Capabilities.}

The architecture allows neural networks to learn algorithmic behaviors such as

\begin{itemize}
		\item graph traversal
		\item path finding
		\item question answering
		\item relational reasoning
\end{itemize}

These tasks require structured memory and sequential retrieval.

\paragraph{Key Results.}

The paper demonstrates strong performance on several reasoning tasks:

\begin{itemize}
		\item family tree inference
		\item graph navigation
		\item synthetic question answering
\end{itemize}

The model learns to store structured representations and retrieve them
in ways resembling classical data structures.

\paragraph{Limitations.}

Despite its flexibility, the DNC architecture has several drawbacks:

\begin{itemize}
		\item Training instability due to complex memory dynamics
		\item Large computational overhead from maintaining link matrices
		\item Difficulty scaling to very large memories
\end{itemize}

As a result, later work explored simpler memory mechanisms such as attention-based
transformers.

\paragraph{Takeaway.}

Differentiable Neural Computers represent an early attempt to combine
neural networks with explicit algorithmic memory. While later architectures
such as transformers replaced many of these mechanisms, DNCs introduced
important ideas about differentiable memory addressing and structured
neural computation.

\end{document}
